% =====================================================================
%  c00.tex —— §1 導入（原书 p.42–43；PDF p.1–p.2 全文 + p.3 顶部两段）
%  编号公式：本文第一条 (1)（原书 p.42）
%  无编号展示：Tr(C) := H^0(Tw(C))
%  定理类环境：无
%  注：跨块引用一律写原书字面节号（第 N 节），文献引用写 [N]，
%      以免单块编译出现 undefined reference / citation。
% =====================================================================

% ===== 原书 p.42 =====
\section{引言}
\label{sec:intro}

作为镜像对称猜想的范畴论表述，Kontsevich 提出的同调镜像对称[14]断言：
辛流形$M$上的 Fukaya 范畴$Fuk(M)$与$M$的镜像对偶复流形$\hat{M}$上相干层构成的
导出范畴$D^{b}(\coh(\hat{M}))$等价。Fukaya 范畴$Fuk(M)$是$\ainf$范畴，
为比较它与导出范畴$D^{b}(\coh(\hat{M}))$，Kontsevich 提出从$\ainf$范畴
$\mathcal{C}$构造三角范畴的方法（[14]；这是 Bondal--Kapranov 的 DG 范畴构造
[2]在$\ainf$范畴情形的推广）。该三角范畴记作$D^{b}(\mathcal{C})$，
常称为$\mathcal{C}$的导出范畴。进而，对$Fuk(M)$与$D^{b}(\coh(\hat{M}))$，
三角范畴的等价
\begin{equation}
  D^{b}(Fuk(M)) \simeq D^{b}(\coh(\hat{M}))
  \label{eq:c00-1}
\end{equation}
成立——这就是同调镜像对称猜想的一种表述[14]。

\transnote{
式~\eqref{eq:c00-1}两侧处在不同的范畴层级：右端是通常意义下相干层的导出范畴，
左端要经本节所述的「单侧扭复形—零次上同调」两步方能得到。原文把二者都称为
「导出范畴」，正是这一构造使比较成为可能。
}

从$\ainf$范畴$\mathcal{C}$构造三角范畴$D^{b}(\mathcal{C})=\Tr(\mathcal{C})$的方法
由以下步骤组成（第 3 节）。

\begin{itemize}[leftmargin=2em,itemsep=2pt,topsep=4pt]
  \item 构造以$\mathcal{C}$的单侧扭复形（one-sided twisted complex）为对象的
        $\ainf$范畴$\Tw(\mathcal{C})$。
  \item 把$\Tr(\mathcal{C})$定义为$\Tw(\mathcal{C})$的零次上同调。
\end{itemize}
\[
  \Tr(\mathcal{C}) := H^{0}(\Tw(\mathcal{C})).
\]

对三角范畴$\mathcal{T}$，若存在$\ainf$范畴$\mathcal{C}$使
$\mathcal{T}\simeq\Tr(\mathcal{C})$，就把$\Tw(\mathcal{C})$称为$\mathcal{T}$的
\textbf{$\ainf$增广}（\mbox{$\ainf$-enhancement}）。此时，可以把$\ainf$范畴
$\mathcal{C}$视为三角范畴$\mathcal{T}$的生成系。一般地，由米田嵌入，可从
$\ainf$范畴构造出与之$\ainf$拟同型（即$\ainf$等价）的 DG 范畴
（cf.\ [4, 20, 11]），因此$\ainf$增广的存在与 DG 增广的存在等价；
但三角范畴未必总有 DG 增广[19]。不过在相干层的导出范畴这类三角范畴的
主要例子上，只要满足一定的良好条件，DG 增广就存在。

% ===== 原书 p.43 =====
$\Tr(\mathcal{C})$成为$\ainf$范畴$\mathcal{C}$的不变量。即，若$\mathcal{C}$、
$\mathcal{C}'$互相$\ainf$等价，则存在三角范畴的等价
$\Tr(\mathcal{C})\simeq\Tr(\mathcal{C}')$（引理 3.2）。因此，给定一组具体的
镜像对偶流形时，证明同调镜像对称\eqref{eq:c00-1}的有力方法之一，就是找出
$D^{b}(Fuk(M))$、$D^{b}(\coh(\hat{M}))$作为生成系的$\ainf$范畴$\mathcal{C}$、
$\mathcal{C}'$，并证明$\ainf$等价$\mathcal{C}\simeq\mathcal{C}'$。

另一方面，作为反方向的问题：从$\ainf$范畴$\mathcal{C}$构造三角范畴
$\Tr(\mathcal{C})$时，信息是否会丢失？实际上，已有这样的例子：互不$\ainf$等价的
$\ainf$范畴$\mathcal{C}$、$\mathcal{C}'$，却存在$\ainf$等价
$\Tw(\mathcal{C})\simeq\Tw(\mathcal{C}')$——只要把通常的范畴$\mathcal{C}$、
$\mathcal{C}'$视为微分与高次$\ainf$积都取零的$\ainf$范畴即可
（cf.\ [9]）。这意味着，作为$\Tw(\mathcal{C})$的生成系，可以存在互不
$\ainf$等价的两个$\ainf$范畴$\mathcal{C}$、$\mathcal{C}'$。事实上，由构造法，
$\ainf$范畴$\Tw(\mathcal{C})$以$\mathcal{C}$为全子$\ainf$范畴；而若
$\Tw(\mathcal{C})\simeq\Tw(\mathcal{C}')$，则与$\mathcal{C}'$等价的$\ainf$范畴
也成为$\Tw(\mathcal{C})$的全子$\ainf$范畴。

\transnote{
这里的「生成系」原文未加限定。全文它在两种意义下使用：第 4 节问题 4.1 中指的是
经典生成系（所生成的三角范畴含全部对象），第 7 节讨论例外生成系时用的是强生成系
（更严）。读到「生成系」一词时，宜按上下文判别。
}

如此，应考虑的问题便是：从$\ainf$范畴$\Tw(\mathcal{C})$得到三角范畴
$\Tr(\mathcal{C})$时，是否会丢失信息？换言之，三角范畴的$\ainf$增广只要存在，
是否唯一（问题 4.1）？

$\Tw(\mathcal{C})$的零次上同调$\Tr(\mathcal{C}):=H^{0}(\Tw(\mathcal{C}))$，
是对构成复形的态射空间取零次上同调，并在这些空间之间放入由$\Tw(\mathcal{C})$的
$\ainf$构造的二次积诱导出的积构造（态射的合成）。因此$\Tw(\mathcal{C})$的高次
$\ainf$积构造全部被舍弃，按通常想法似乎丢失了相当多的信息。然而实际上，
$\mathcal{C}$的高次积构造的信息有相当一部分反映在$\Tr(\mathcal{C})$的积构造中
而留存下来。作为$\Tr(\mathcal{C})$的积构造的具体表示，第 5 节将说明：
由$\mathcal{C}$定义的扭积（twisted product），一般包含分次微分代数，或者分次
微分范畴$\mathcal{C}$的 Massey 积及其各种推广。

然而归根结底，$\mathcal{C}$的扭积并不能完全复原$\mathcal{C}$的$\ainf$构造。
事实上，存在互不$\ainf$等价的$\ainf$范畴$\mathcal{C}$、$\mathcal{C}'$，使三角范畴
的等价$\Tr(\mathcal{C})\simeq\Tr(\mathcal{C}')$成立，但它不能提升为$\ainf$等价
$\Tw(\mathcal{C})\simeq\Tw(\mathcal{C}')$；这正是$\ainf$增广唯一性问题（问题 4.1）
的反例[12]。介绍这方面内容（第 6 节）是本文的主要目的。

在[16]中，Lunts 与 Orlov 对相当广泛的一类三角范畴证明了 DG 增广的
唯一性。但那里的 DG 增广唯一性并非$\ainf$等价类的唯一性，而是关于由拟函子
（quasi-functor）这一概念定义的等价类的唯一性，因此与本文的讨论并不矛盾。

\transnote{
这句是全文最易被误读的地方：拟函子的等价类比$\ainf$等价类粗得多——两个 DG 范畴
可以在拟函子意义下等价，却不在$\ainf$等价意义下等价。因此 Lunts--Orlov 的唯一性
结论并不触及本文所问的$\ainf$增广唯一性，两者并不冲突。
}

以下兼作记号说明：第 2 节复习$\ainf$范畴的基本定义与性质，第 3 节复习从$\ainf$
范畴构造三角范畴的方法，均只取最低限度。在这些准备之下，第 4 节陈述本文的
主问题，第 5 节定义扭积，随后第 6 节构造一个$\ainf$增广不唯一的三角范畴的例子。
最后第 7 节讨论一类能简单看出$\ainf$增广唯一的三角范畴。

% ===== 原书 p.44（§2 标题之前的两段，属 §1） =====
本文中，提到（$\mathbb{Z}$分次）向量空间时，总指某个固定特征零的域$K$上的
向量空间；三角范畴、$\ainf$范畴也都在$K$上考虑。
